% !PRINTED 233
\weberpage{233}

Da sich die linke Seite durch Vertauschung von $\sigma$ mit $-\sigma$ nicht \"andert, so ist es eine reine quadratische Gleichung und sie giebt mit R\"ucksicht auf \S.\,73, (12)
\[
B_{0,3}^2\,\sigma^2 - 5\,D = 0,
\]
wenn $D$ die Discriminante der gegebenen Gleichung ist. Wir bekommen also nach (2)
\weq{12}{5\,a_0\,a_5\,\sigma = \sqrt{5\,D},}
worin das Vorzeichen beliebig ist.

\begin{center}
\S.\ 75.

\gesperrt{Normalform der Gleichung f\"unften Grades.}
\end{center}

Wie schon fr\"uher bemerkt, ist das haupts\"achlichste Ziel dieser Betrachtungen, eine Normalform der Gleichung f\"unften Grades herzustellen, die nur von \gesperrt{einem} unbestimmten Co\"efficienten, \gesperrt{einem Parameter}, abh\"angt. Eine solche Normalform ist die \gesperrt{Bring-Jerrard'sche} Form. Um diese zu erhalten, haben wir nach (4) des vorigen Paragraphen
\weq{1}{y = t_3\,F_0 + t_2\,(\alpha_2\,F_1 + \alpha_1\,F_2 + \alpha_0\,F_3)}
zu setzen, so dass identisch
\[
S(y) = 0, \quad S(y^2) = 0
\]
wird, und dann ist das Verh\"altniss $t_3 : t_2$ aus der cubischen Gleichung
\[
S(y^3) = 0
\]
zu bilden. Diese cubische Gleichung l\"asst sich wirklich bilden, wenn auch ihr Ausdruck lang wird. Die hierzu n\"othigen Formeln sind von \gesperrt{Cayley} berechnet\textsuperscript{1)}.

Es ist das Ergebniss einer merkw\"urdigen Untersuchung von \gesperrt{Gordan}, dass man eine andere, die \gesperrt{Brioschi'sche} Normalform ohne neue Irrationalit\"at erhalten kann, und wir wollen zum Beschluss dieser Betrachtungen \"uber die Tschirnhausen-Transformation dies Resultat noch ableiten\textsuperscript{2)}.

Wir halten an den Voraussetzungen des vorigen Paragraphen fest und setzen f\"ur den Augenblick zur Abk\"urzung
%CONT>>

\bigskip
\noindent\rule{4cm}{0.4pt}

\noindent\textsuperscript{1)} \gesperrt{Cayley}, on Tschirnhausen's Transformation, Phil. Trans. 1861, Mathematical Papers, Tom. IV, Nr. 275.

\noindent\textsuperscript{2)} \gesperrt{Gordan}, Mathematische Annalen, Bd. 28, 1886.
